% Options for packages loaded elsewhere
\PassOptionsToPackage{unicode}{hyperref}
\PassOptionsToPackage{hyphens}{url}
\documentclass[
  12pt,
]{ctexart}
\usepackage{xcolor}
\usepackage{amsmath,amssymb}
\setcounter{secnumdepth}{-\maxdimen} % remove section numbering
\usepackage{iftex}
\ifPDFTeX
  \usepackage[T1]{fontenc}
  \usepackage[utf8]{inputenc}
  \usepackage{textcomp} % provide euro and other symbols
\else % if luatex or xetex
  \usepackage{unicode-math} % this also loads fontspec
  \defaultfontfeatures{Scale=MatchLowercase}
  \defaultfontfeatures[\rmfamily]{Ligatures=TeX,Scale=1}
\fi
\usepackage{lmodern}
\ifPDFTeX\else
  % xetex/luatex font selection
\fi
% Use upquote if available, for straight quotes in verbatim environments
\IfFileExists{upquote.sty}{\usepackage{upquote}}{}
\IfFileExists{microtype.sty}{% use microtype if available
  \usepackage[]{microtype}
  \UseMicrotypeSet[protrusion]{basicmath} % disable protrusion for tt fonts
}{}
\makeatletter
\@ifundefined{KOMAClassName}{% if non-KOMA class
  \IfFileExists{parskip.sty}{%
    \usepackage{parskip}
  }{% else
    \setlength{\parindent}{0pt}
    \setlength{\parskip}{6pt plus 2pt minus 1pt}}
}{% if KOMA class
  \KOMAoptions{parskip=half}}
\makeatother
\usepackage{longtable,booktabs,array}
\usepackage{calc} % for calculating minipage widths
% Correct order of tables after \paragraph or \subparagraph
\usepackage{etoolbox}
\makeatletter
\patchcmd\longtable{\par}{\if@noskipsec\mbox{}\fi\par}{}{}
\makeatother
% Allow footnotes in longtable head/foot
\IfFileExists{footnotehyper.sty}{\usepackage{footnotehyper}}{\usepackage{footnote}}
\makesavenoteenv{longtable}
\ifLuaTeX
  \usepackage{luacolor}
  \usepackage[soul]{lua-ul}
\else
  \usepackage{soul}
\fi
\setlength{\emergencystretch}{3em} % prevent overfull lines
\providecommand{\tightlist}{%
  \setlength{\itemsep}{0pt}\setlength{\parskip}{0pt}}
\usepackage{bookmark}
\IfFileExists{xurl.sty}{\usepackage{xurl}}{} % add URL line breaks if available
\urlstyle{same}
\hypersetup{
  hidelinks,
  pdfcreator={LaTeX via pandoc}}

\author{SCX}
\date{2026}

\begin{document}

\section{SCX 在 LLM
大模型方向的战略与发展路线图}\label{scx-ux5728-llm-ux5927ux6a21ux578bux65b9ux5411ux7684ux6218ux7565ux4e0eux53d1ux5c55ux8defux7ebfux56fe}

\begin{quote}
基于 2026-06-26 与 GPT 的深度讨论整理 整理日期：2026-06-26
\end{quote}

\begin{center}\rule{0.5\linewidth}{0.5pt}\end{center}

\subsection{一、LLM 对 SCX
为什么重要}\label{ux4e00llm-ux5bf9-scx-ux4e3aux4ec0ux4e48ux91cdux8981}

\subsubsection{1.1 GPT 的核心论证：LLM 数据筛选是 SCX
的''远期外延''}\label{gpt-ux7684ux6838ux5fc3ux8bbaux8bc1llm-ux6570ux636eux7b5bux9009ux662f-scx-ux7684ux8fdcux671fux5916ux5ef6}

GPT 多次强调一个核心判断：

\begin{quote}
\textbf{SCX 在 ACE/MLIP 上已经有了强 prototype，这给了你很强的动机去往
LLM 扩展，但 ACE 成功不等于 LLM 会自动成功。}
\end{quote}

ACE/MLIP 成功的原因是： - 有明确的 DFT oracle（gold label） - 物理
descriptor 状态空间天然有意义 - 冗余非常明显（AIMD 相邻帧相似） -
物理有约束（能量守恒、连续性）

而 LLM 场景会遇到四个新难题： 1.
\textbf{表示空间不等于价值空间}：embedding
相近的两条文本，训练价值可能完全不同 2. \textbf{误差难以定义}：LLM
数据没有绝对 gold label，只能用代理量（loss、perplexity、model
disagreement 等） 3.
\textbf{数据价值不是局部可加的}：一条数据的价值依赖已有数据集和当前模型
4. \textbf{冗余不能只靠相似度判断}：LLM
中''看起来重复''的数据可能增强鲁棒性

\subsubsection{1.2 市场大小对比：材料 vs
LLM}\label{ux5e02ux573aux5927ux5c0fux5bf9ux6bd4ux6750ux6599-vs-llm}

\begin{longtable}[]{@{}
  >{\raggedright\arraybackslash}p{(\linewidth - 4\tabcolsep) * \real{0.2222}}
  >{\raggedright\arraybackslash}p{(\linewidth - 4\tabcolsep) * \real{0.4074}}
  >{\raggedright\arraybackslash}p{(\linewidth - 4\tabcolsep) * \real{0.3704}}@{}}
\toprule\noalign{}
\begin{minipage}[b]{\linewidth}\raggedright
维度
\end{minipage} & \begin{minipage}[b]{\linewidth}\raggedright
材料/MLIP
\end{minipage} & \begin{minipage}[b]{\linewidth}\raggedright
LLM 数据
\end{minipage} \\
\midrule\noalign{}
\endhead
\bottomrule\noalign{}
\endlastfoot
市场规模 & 较小（学术+材料公司） & 巨大（所有 AI 公司） \\
竞争强度 & 低（SCX 是原创方向） & 极高（大厂+学术都在做） \\
SCX 优势 & 强（有 DFT oracle + 物理描述符） &
弱（标签模糊、误差不可定义） \\
商业化难度 & 中等（需要行业合作） & 极高（大厂内部化） \\
\end{longtable}

GPT 的最终判断：\textbf{SCX 不适合正面切入通用 LLM
数据评价，但适合作为远期外延和论文 discussion。}

\subsubsection{1.3 SCX 在 LLM
数据筛选中的独特优势}\label{scx-ux5728-llm-ux6570ux636eux7b5bux9009ux4e2dux7684ux72ecux7279ux4f18ux52bf}

SCX 在 LLM 场景的合理定位不是''又一个更准的 judge''，而是：

\begin{quote}
\textbf{状态条件的 judge 可靠性校准系统 / 评价器编排框架。}
\end{quote}

具体来说，SCX 可以补的是： 1. \textbf{不同 judge
在不同状态下的可靠性评估}：强 GPT 也不是所有状态都强，数学强不一定医学强
2. \textbf{高分数据的真实边际价值判断}：GPT
喜欢''漂亮但重复''的数据，SCX 加入冗余项 3.
\textbf{低分数据的训练价值识别}：覆盖模型薄弱状态的数据可能更值钱 4.
\textbf{评价器成本感知调度}：不同状态切换使用不同评价器，控制成本

\begin{center}\rule{0.5\linewidth}{0.5pt}\end{center}

\subsection{二、LLM 实验需求}\label{ux4e8cllm-ux5b9eux9a8cux9700ux6c42}

\subsubsection{2.1 为什么要做 LLM
实验}\label{ux4e3aux4ec0ux4e48ux8981ux505a-llm-ux5b9eux9a8c}

要说服审稿人和客户，SCX 在 LLM 数据筛选上有效，必须有一个最小可行实验。

但 GPT 强调：\textbf{LLM 实验不能作为 SCX 的核心实验主战场。} LLM
只能作为 discussion/extension，核心实验应该在 ACE/MLIP 和医学图像上。

\subsubsection{2.2 最小可行 LLM
demo}\label{ux6700ux5c0fux53efux884c-llm-demo}

\begin{longtable}[]{@{}ll@{}}
\toprule\noalign{}
项目 & 规格 \\
\midrule\noalign{}
\endhead
\bottomrule\noalign{}
\endlastfoot
模型 & 0.5B-3B 小模型 + LoRA 微调 \\
数据 & 公开 instruction dataset（几千到几万条） \\
硬件 & 单张 4090 足够 \\
\end{longtable}

\textbf{SCX 表示}：sentence embedding 或模型 hidden states

\textbf{误差代理}：初始模型 loss 或多模型 disagreement

\subsubsection{2.3 实验设计}\label{ux5b9eux9a8cux8bbeux8ba1}

\textbf{冗余判断公式：}

\begin{verbatim}
Redundancy(s) = [1 - mean_r(s)] * density(s) * similarity(s)
\end{verbatim}

\textbf{数据压缩策略：} - 保留高误差高密度状态 - 保留边界状态 -
每个冗余状态保留 medoid - 设置随机/多样性对照组

\textbf{验证对比：} - full data（100\%） - random 20\% - diversity 20\%
- high-loss 20\% - SCX 20\%

\textbf{目标：} SCX 20\% 接近 full data 效果，或优于 random/diversity

\subsubsection{2.4
具体技术路线}\label{ux5177ux4f53ux6280ux672fux8defux7ebf}

\textbf{实验 A：冗余压缩} 1. 训练基础模型，提取 embedding 2. 计算 loss
作为残差 3. 聚类形成状态 4. 对每个状态计算冗余分数 5.
只保留代表样本，加权训练 6. 对比 full data / random / uncertainty /
coreset

\textbf{实验 B：高误差不等于高价值} 1. 人为加入 label noise 2. 比较
high-loss sampling、uncertainty sampling、SCX state-wise acquisition 3.
证明 SCX 优先选高密度、一致、可学习的高误差状态（而非噪声点）

\textbf{实验 C：专家路由} 1. 训练多个不同能力的 judge/expert 2. 估计每个
expert 在不同状态上的可靠性 3.
路由决策：为每个样本选择最可靠且成本最优的 judge

\begin{center}\rule{0.5\linewidth}{0.5pt}\end{center}

\subsection{三、发展路线图（时间线）}\label{ux4e09ux53d1ux5c55ux8defux7ebfux56feux65f6ux95f4ux7ebf}

\subsubsection{阶段 1：当前（0-2 月）------
立住根}\label{ux9636ux6bb5-1ux5f53ux524d0-2-ux6708-ux7acbux4f4fux6839}

\textbf{关键里程碑：} - {[} {]} 完成 EGP Paper 1（ACE gauge-normalized
expert merging） - {[} {]} 新建私人 SCX 仓库，clean-room 重写代码 - {[}
{]} 完成 DEVELOPMENT\_LOG.md，建立证据隔离 - {[} {]} 查 IP
协议，做专利评估 - {[} {]} 用 MedMNIST 跑通第一个 SCX-Compress demo -
{[} {]} 用最小 instruction dataset 跑通第一个 LLM
小实验（验证趋势，不追求大规模）

\textbf{目标：} 证明 SCX 不是 PPT，在 ACE/MLIP 上有强原型

\subsubsection{阶段 2：1-2 月 ------
建理论地基}\label{ux9636ux6bb5-21-2-ux6708-ux5efaux7406ux8bbaux5730ux57fa}

\textbf{关键里程碑：} - {[} {]} arXiv 发布 SCX-Theory（数学定义 +
可识别性边界 + 合成实验） - {[} {]} 抢占命名权：State-Conditioned
eXpertise - {[} {]} 完成 SCX-Core
开源（轻量版，做到''足以复现论文，不足以替代产品''） - {[} {]} 跑通
HAM10000/CheXpert 医学图像实验 - {[} {]} 完成 LLM 小规模 SFT
验证实验（证明 SCX 选择的 20\% 数据趋势有效） - {[} {]} 校外资 IP
律师咨询

\textbf{目标：} 建立 SCX 的概念占位和理论地基，LLM 作为
discussion/extension

\subsubsection{阶段 3：3-6 月 ------
攻大论文}\label{ux9636ux6bb5-33-6-ux6708-ux653bux5927ux8bbaux6587}

\textbf{关键里程碑：} - {[} {]} 完成 SCX-Sim 大论文（DFT/MLIP + FEM/PDE
+ 医学图像 + OPC/TCAD toy） - {[} {]} 目标期刊：Nature Computational
Science / Nature Communications - {[} {]} 找 1-2 个甲方做 paid pilot -
{[} {]} 完成 SCX-Health 开源框架 - {[} {]} 构建 SCX Potential Compiler
原型（ACE + NEP + MACE + DeepMD 多专家蒸馏） - {[} {]} 完善 LLM
实验：在更大 SFT 数据集上验证 SCX routing 的性价比

\textbf{目标：} 证明 SCX 是跨领域的科学/工程仿真数据价值评估框架；LLM
数据评价的论文讨论充分展开

\subsubsection{阶段 4：6-12 月 ------
商业化落地}\label{ux9636ux6bb5-46-12-ux6708-ux5546ux4e1aux5316ux843dux5730}

\textbf{关键里程碑：} - {[} {]} 成立 SCX Working Group /
Consortium（多公司参与） - {[} {]} 建立 SCX 认证体系雏形（SCX-certified
dataset 等） - {[} {]} 找 CEO/COO，自己退到 scientific founder 角色 -
{[} {]} 将 LLM 方向的 SCX-Judge 作为企业版功能模块（本地部署 +
数据审计报告） - {[} {]}
正式接触华为（健康/昇腾生态）、小米（可穿戴/IoT）、半导体（国产
TCAD/CMP） - {[} {]} 将 ``SCX 评价器编排系统''
包装为企业可购买的数据审计服务

\textbf{目标：} 从学术论文走向产业标准

\subsubsection{LLM
方向在整个路线图中的位置}\label{llm-ux65b9ux5411ux5728ux6574ux4e2aux8defux7ebfux56feux4e2dux7684ux4f4dux7f6e}

\begin{verbatim}
阶段 1: LLM 小实验（验证趋势）
      ↓
阶段 2: LLM 嵌入 SCX-Theory（作为讨论/extension）
      ↓
阶段 3: LLM 实验完善（作为 SCX-Sim 的一个外延验证）
      ↓
阶段 4: SCX-Judge 模块商业化（企业版功能）
\end{verbatim}

\textbf{核心原则：LLM 是远期外延，不是主战场。}

\begin{center}\rule{0.5\linewidth}{0.5pt}\end{center}

\subsection{四、资源配置}\label{ux56dbux8d44ux6e90ux914dux7f6e}

\subsubsection{4.1 人力需求}\label{ux4ebaux529bux9700ux6c42}

\begin{longtable}[]{@{}lll@{}}
\toprule\noalign{}
角色 & 必要性 & 来源建议 \\
\midrule\noalign{}
\endhead
\bottomrule\noalign{}
\endlastfoot
你自己（学术/算法核心） & 必须 & SCX 理论+核心代码 \\
1-2 名学生/合作者 & 需要 & 协助实验和代码 \\
运营负责人 & 阶段 2-3 需要 & 管合同、客户、项目排期 \\
工程负责人 & 阶段 3-4 需要 & 管平台、部署、CLI 工具 \\
法务/IP 顾问 & 优先找 & 管专利、许可、CLA、合作协议 \\
商务顾问/产业导师 & 阶段 3 需要 & 帮你找甲方，不被甲方牵着走 \\
\end{longtable}

\subsubsection{4.2 算力需求}\label{ux7b97ux529bux9700ux6c42}

\begin{longtable}[]{@{}lll@{}}
\toprule\noalign{}
实验类型 & 最低配置 & 推荐配置 \\
\midrule\noalign{}
\endhead
\bottomrule\noalign{}
\endlastfoot
MedMNIST / 小型图像实验 & 笔记本 CPU & 4090 \\
ACE/MLIP 势函数 & 单卡 4090 & 单卡 4090 \\
LLM 最小实验（0.5B-3B + LoRA） & 单卡 4090（24GB） & 单卡 4090 \\
更大规模 LLM SFT & 单卡 4090（可跑 7B 量级） & A100/多卡 \\
SCX-Potential Compiler & 单卡 4090 & 多卡并行 \\
工业级大规模 pipeline & - & 需要集群 \\
\end{longtable}

\textbf{结论：4090 可以跑通全部学术版 demo，包括 LLM
最小实验。短期不需要追加算力投资。}

\subsubsection{4.3
资金需求估算}\label{ux8d44ux91d1ux9700ux6c42ux4f30ux7b97}

\begin{longtable}[]{@{}
  >{\raggedright\arraybackslash}p{(\linewidth - 4\tabcolsep) * \real{0.2308}}
  >{\raggedright\arraybackslash}p{(\linewidth - 4\tabcolsep) * \real{0.3846}}
  >{\raggedright\arraybackslash}p{(\linewidth - 4\tabcolsep) * \real{0.3846}}@{}}
\toprule\noalign{}
\begin{minipage}[b]{\linewidth}\raggedright
阶段
\end{minipage} & \begin{minipage}[b]{\linewidth}\raggedright
主要支出
\end{minipage} & \begin{minipage}[b]{\linewidth}\raggedright
预算估算
\end{minipage} \\
\midrule\noalign{}
\endhead
\bottomrule\noalign{}
\endlastfoot
阶段 1（0-2月） & 校外 IP 律师（1次咨询）、论文投稿费 &
\textasciitilde0.5-1 万 \\
阶段 2（1-2月） & arXiv 费用、开源仓库托管、律师持续咨询 &
\textasciitilde1-2 万 \\
阶段 3（3-6月） & Nature 系列 OA 版面费（\textasciitilde2-4
万人民币）、vaspkit/软件许可、合作差旅 & \textasciitilde5-10 万 \\
阶段 4（6-12月） & 运营/工程人力、AWS/超算、法务、商务、公司注册 &
\textasciitilde20-50 万 \\
\end{longtable}

\textbf{资金来源建议：} 1. 阶段 1-2：个人资金 +
课题组经费（如果切割清楚） 2. 阶段 3：paid pilot 收入 + 赞助研究 3. 阶段
4：membership fee + 企业版授权 + 可能的小规模种子轮

\begin{center}\rule{0.5\linewidth}{0.5pt}\end{center}

\subsection{五、优先级决策}\label{ux4e94ux4f18ux5148ux7ea7ux51b3ux7b56}

\subsubsection{5.1
现在必须做（高优先级）}\label{ux73b0ux5728ux5fc5ux987bux505aux9ad8ux4f18ux5148ux7ea7}

\begin{enumerate}
\def\labelenumi{\arabic{enumi}.}
\tightlist
\item
  \textbf{EGP Paper 1 (ACE gauge merging)}：这是你已有的最强资产，优先出
\item
  \textbf{证据隔离}：新建私人仓库、clean-room 写代码、写
  DEVELOPMENT\_LOG.md
\item
  \textbf{IP 保护}：查协议、问律师、考虑专利/软著/时间戳
\item
  \textbf{SCX-Core 最小实现}：核心数学 + 合成数据 demo
\item
  \textbf{MedMNIST 跑通}：用公开医学数据做第一个外部验证
\item
  \textbf{最小 LLM 实验}：用 0.5B-3B 模型 + LoRA 验证趋势
\end{enumerate}

\subsubsection{5.2
可以等论文发表后}\label{ux53efux4ee5ux7b49ux8bbaux6587ux53d1ux8868ux540e}

\begin{itemize}
\tightlist
\item
  完整开源 SCX-Health
\item
  找甲方做 paid pilot
\item
  建立产业联盟
\item
  正式公司化
\item
  大规模 LLM 数据评价实验
\end{itemize}

\subsubsection{5.3
需要合作者/资金后}\label{ux9700ux8981ux5408ux4f5cux8005ux8d44ux91d1ux540e}

\begin{itemize}
\tightlist
\item
  TCAD/OPC/CMP 工业数据实验（需要行业合作）
\item
  多中心临床医学验证（需要医院合作）
\item
  华为/小米可穿戴数据试点（需要企业协议）
\item
  分布式大规模 SCX 计算平台（需要算力投资）
\end{itemize}

\subsubsection{5.4 对 LLM
的精力分配建议}\label{ux5bf9-llm-ux7684ux7cbeux529bux5206ux914dux5efaux8bae}

GPT 的原始建议值得全文引用：

\begin{quote}
\textbf{``大模型不适合做你的主战场。大模型是远期外延，不是核心根据地。''}

\textbf{``大模型数据评价已经被大厂内部化；你一个博士硬冲很难。''}

\textbf{``SCX 在 LLM
上的合理切口：不是造主任医师，而是造医院的分诊、会诊、质控和病历审计系统。''}
\end{quote}

具体来说：

\begin{longtable}[]{@{}
  >{\raggedright\arraybackslash}p{(\linewidth - 4\tabcolsep) * \real{0.3448}}
  >{\raggedright\arraybackslash}p{(\linewidth - 4\tabcolsep) * \real{0.4483}}
  >{\raggedright\arraybackslash}p{(\linewidth - 4\tabcolsep) * \real{0.2069}}@{}}
\toprule\noalign{}
\begin{minipage}[b]{\linewidth}\raggedright
LLM 方向
\end{minipage} & \begin{minipage}[b]{\linewidth}\raggedright
该花多少精力
\end{minipage} & \begin{minipage}[b]{\linewidth}\raggedright
原因
\end{minipage} \\
\midrule\noalign{}
\endhead
\bottomrule\noalign{}
\endlastfoot
SCX-Theory 中讨论 LLM 作为外延 & 5\% & 论文需要提及，但不展开 \\
最小 LLM 实验验证趋势 & 10\% & 有比没有好，但不作为核心证据 \\
LLM-as-Judge 可靠性校准论文 & 15\% & 如果有余力可以单独发短文 \\
通用 LLM 数据评价平台 & 不投入 & 大厂内部在做，很难差异化 \\
垂直领域 LLM 数据审计（法律/医学/科学） & 20\% &
大厂不愿意细做的领域，有空间 \\
\end{longtable}

\textbf{最终结论：把 70\% 精力放在 SCX-MLIP + SCX-Theory + SCX-Sim
三条主线上，20\% 放在医学/视觉开源，10\% 放在 LLM 外延验证。}

\begin{center}\rule{0.5\linewidth}{0.5pt}\end{center}

\subsection{六、SCX 在 LLM
的竞争分析}\label{ux516dscx-ux5728-llm-ux7684ux7adeux4e89ux5206ux6790}

\subsubsection{6.1 现有 9
类主流方法}\label{ux73b0ux6709-9-ux7c7bux4e3bux6d41ux65b9ux6cd5}

\begin{longtable}[]{@{}
  >{\raggedright\arraybackslash}p{(\linewidth - 4\tabcolsep) * \real{0.2143}}
  >{\raggedright\arraybackslash}p{(\linewidth - 4\tabcolsep) * \real{0.3571}}
  >{\raggedright\arraybackslash}p{(\linewidth - 4\tabcolsep) * \real{0.4286}}@{}}
\toprule\noalign{}
\begin{minipage}[b]{\linewidth}\raggedright
方法
\end{minipage} & \begin{minipage}[b]{\linewidth}\raggedright
核心思想
\end{minipage} & \begin{minipage}[b]{\linewidth}\raggedright
SCX 的关系
\end{minipage} \\
\midrule\noalign{}
\endhead
\bottomrule\noalign{}
\endlastfoot
规则清洗（FineWeb） & 启发式过滤垃圾 & SCX 可以补充状态条件质量判断 \\
质量分类器（FinerWeb） & 强模型标一小批，训练轻量分类器 & SCX
可加''状态条件''：Q(x \\
Perplexity/Loss（Rho-1） & 用 loss 判断 token 价值 & SCX 可上升到
state-level selection \\
目标分布匹配（DSIR） & 选最像目标分布的数据 & SCX 可扩展为状态条件
DSIR \\
Influence/Gradient（LoGra/LESS） & 看数据对目标 loss 的贡献 &
计算贵，SCX 可在高价值状态上调用 \\
LLM-as-a-Judge & 用强模型评价 & SCX 做 judge 可靠性校准 \\
多样性/去重 & 避免重复数据 & SCX 集成到价值公式 V(x)=Q*(1-D) \\
可验证任务（数学/代码） & 用标准答案验证 & SCX 设为高可靠专家 \\
Benchmark-driven（DataComp-LM） & 最终看训练效果 & SCX 的闭环验证 \\
\end{longtable}

\subsubsection{6.2 SCX
的框架优势}\label{scx-ux7684ux6846ux67b6ux4f18ux52bf}

现有方法大多是\textbf{单信号或单阶段}，而 SCX 回答的是：

\begin{quote}
\textbf{在当前状态 (s)、当前模型 (M)、当前数据集 (D)、可用评价器集合 (J)
下，这个样本应该采取什么动作？}
\end{quote}

动作不是只有 keep/drop，而是：

\begin{verbatim}
keep | drop | downweight | deduplicate | relabel
| route to stronger judge | verify with tool
| use for SFT | use for preference data | use for eval
\end{verbatim}

核心公式：

\begin{verbatim}
Q_SCX(x) = sum_j w_j(s(x)) * Q_j(x)

其中 w_j(s) ∝ exp(alpha * A_j(s) - lambda * C_j)

- Q_j(x): 第 j 个评价器的分数
- A_j(s): 该评价器在状态 s 上的可靠性
- C_j: 评价成本
\end{verbatim}

\subsubsection{6.3 不要正面硬刚的
Statement}\label{ux4e0dux8981ux6b63ux9762ux786cux521aux7684-statement}

\begin{quote}
\st{我提出一个更准的数据质量算法。}
\end{quote}

应该说： \textgreater{}
\textbf{我提出一个状态条件的数据质量评价编排框架。它把 LLM judge、reward
model、verifier、influence score、loss
dynamics、dedup、target-distribution matching
统一到同一个决策系统里，并学习每个评价器在不同数据状态下的可靠性。}

\begin{center}\rule{0.5\linewidth}{0.5pt}\end{center}

\subsection{七、风险与注意事项}\label{ux4e03ux98ceux9669ux4e0eux6ce8ux610fux4e8bux9879}

\subsubsection{7.1 LLM
方向的主要风险}\label{llm-ux65b9ux5411ux7684ux4e3bux8981ux98ceux9669}

\begin{enumerate}
\def\labelenumi{\arabic{enumi}.}
\tightlist
\item
  \textbf{被大厂平替的风险}：通用 LLM 数据评价大厂都在做，差异化很难
\item
  \textbf{证明成本高}：LLM
  数据价值需要大规模训练实验验证，小实验说服力有限
\item
  \textbf{标签模糊}：没有 gold label，SCX 只能输出弱判断（noise-risk
  而非 noise）
\item
  \textbf{Goodhart 问题}：企业可能优化数据迎合 SCX 评分而非真正提升模型
\end{enumerate}

\subsubsection{7.2
与其他方向的资源冲突}\label{ux4e0eux5176ux4ed6ux65b9ux5411ux7684ux8d44ux6e90ux51b2ux7a81}

LLM 方向投入过多会影响主战场（MLIP、FEM、医学）。\textbf{必须控制 LLM
方向不超过总精力的 10-15\%。}

\subsubsection{7.3 关键建议}\label{ux5173ux952eux5efaux8bae}

\begin{enumerate}
\def\labelenumi{\arabic{enumi}.}
\tightlist
\item
  \textbf{不要为了 LLM 放弃 ACE/MLIP}------那是你最有原创证据链的地方
\item
  \textbf{不要声称''替代大厂 judge''}------要说''补充和校准''
\item
  \textbf{LLM 实验只做趋势验证}------不要追求大规模，小数据集跑通即可
\item
  \textbf{垂直领域优先}------法律、材料科学、医学问答等
  SCX-MedJudge/SCX-STEMJudge 比通用 LLM judge 更有空间
\item
  \textbf{先有根据地，再谈外延}------SCX-MLIP 和 SCX-Theory 出不来，LLM
  方向没有根基
\end{enumerate}

\begin{center}\rule{0.5\linewidth}{0.5pt}\end{center}

\subsection{八、核心行动清单（Next
Actions）}\label{ux516bux6838ux5fc3ux884cux52a8ux6e05ux5355next-actions}

\subsubsection{本周}\label{ux672cux5468}

\begin{itemize}
\tightlist
\item[$\square$]
  新建私人 SCX 仓库（个人邮箱、个人电脑）
\item[$\square$]
  写 DEVELOPMENT\_LOG.md 第一版
\item[$\square$]
  查入学/助研协议和学校 IP 政策
\item[$\square$]
  定 EGP Paper 1 论文框架
\end{itemize}

\subsubsection{本月}\label{ux672cux6708}

\begin{itemize}
\tightlist
\item[$\square$]
  clean-room 写 SCX-Core 代码
\item[$\square$]
  MedMNIST 跑通第一个 SCX-Compress 实验
\item[$\square$]
  用 0.5B-3B 模型 + LoRA 跑通最小 LLM 实验
\item[$\square$]
  联系校外 IP 律师
\end{itemize}

\subsubsection{下月}\label{ux4e0bux6708}

\begin{itemize}
\tightlist
\item[$\square$]
  arXiv SCX-Theory 草稿
\item[$\square$]
  开源 SCX-Core（轻量版）
\item[$\square$]
  完成 HAM10000 实验
\item[$\square$]
  论文投出 EGP Paper 1
\end{itemize}

\subsubsection{季度目标}\label{ux5b63ux5ea6ux76eeux6807}

\begin{itemize}
\tightlist
\item[$\square$]
  SCX-Sim 大论文成型
\item[$\square$]
  跑通 FEM/PDE toy demo
\item[$\square$]
  完成 SCX-Health 开源框架
\item[$\square$]
  LLM 方向论文讨论/extension 写完
\item[$\square$]
  开始接触早期甲方（材料/医学方向优先）
\end{itemize}

\begin{center}\rule{0.5\linewidth}{0.5pt}\end{center}

\begin{quote}
\textbf{一句话总结：SCX 在 LLM 方向有框架优势，但不应作为主战场。用 10\%
精力做 LLM 外延验证，70\% 精力做 MLIP+Theory+Sim 主线，20\%
做医学开源。先有根据地，再谈大模型。}
\end{quote}

\end{document}
